%==============================================================================
% 2. RELATED WORK
%==============================================================================
% Drafted from paper/.research/landscape.md (refreshed 2026-05-29) and design.md.
% Cite keys used (all defined in refs.bib):
%   gbrain, zep_graphiti_2025, mem0_2024, graphify_oss            (systems)
%   maharana2024locomo, wu2024longmemeval, beam2025,
%   hu2025memoryagentbench, memorybench2025                        (benchmarks)
%   gebru2021datasheets, croissant2024   (D&B methodology -- now cited in design.tex
%                                          / main.tex; methodology subsection relocated)
% Scope discipline: the public field evaluated here is EXACTLY four systems
% (gbrain, zep-cloud, mem0-platform, graphify-oss). No authors'-own system is
% referenced anywhere.
%==============================================================================
\section{Related Work}
\label{sec:related}

\bench{} sits at the intersection of three lines of work: the agent-memory
systems it evaluates, the benchmarks that currently measure them, and the
broader literature on long-horizon and temporal reasoning. We close the section
by situating \bench{} against existing benchmarks on the axes that distinguish
organizational memory (Table~\ref{tab:bench-comparison}).

\subsection{AI Agent Memory Systems}
\label{sec:related:systems}

A memory system gives an agent persistent state beyond its context window: it
ingests a stream of text, decides what to store and in what representation, and
retrieves the relevant fragments when a later query arrives. The four publicly
available systems we evaluate span the dominant architectural families and make
sharply different design bets, which is precisely why a benchmark that exercises
them under a single harness is informative.

\textbf{gbrain} is an open hybrid-retrieval system~\citep{gbrain} that combines
lexical (BM25), dense vector, and graph-based retrieval over an extracted store,
then fuses the three signals at query time. The hybrid design is intended to
recover both surface lexical matches and semantically paraphrased evidence while
using graph structure to follow relations between entities.

\textbf{Zep}, built on the Graphiti engine~\citep{zep_graphiti_2025}, maintains a
\emph{temporal knowledge graph}: ingested text is distilled into entities and
relations whose edges carry validity intervals, so the graph can in principle be
queried for the state of the world as of a given date and can represent facts that
were later invalidated. Of the four systems it is architecturally closest to the
bi-temporal structure \bench{} probes, which makes its performance a useful test
of whether a temporal-graph representation, on its own, suffices for
organizational-scale questions.

\textbf{Mem0}~\citep{mem0_2024} follows an extract-and-store paradigm aimed at
production deployment: an extraction step converts incoming messages into compact
natural-language memories, with subsequent operations to add, update, or
invalidate them as new information arrives. The representation is optimized for
scalable low-latency retrieval of salient personal facts rather than for
reconstructing a multi-author audit trail.

\textbf{graphify}~\citep{graphify_oss} is an open-source, embedding-free library
that stores memories as a graph and retrieves by explicit BFS/DFS traversal rather
than by vector similarity. By avoiding dense embeddings it removes one source of
opaque approximation, but it also forgoes the soft semantic matching the other
systems rely on, trading recall on paraphrased evidence for transparent, traversal-based
retrieval.

{\sloppy These systems are evaluated head-to-head under a uniform protocol, each
on its vendor-recommended configuration, in Section~\ref{sec:experiments}.\par}

\subsection{Existing Memory Benchmarks}
\label{sec:related:benchmarks}

The benchmarks that currently define progress on agent memory share one structural
assumption: memory is a \emph{single narrator's} history, possibly very long.

\textbf{LoCoMo}~\citep{maharana2024locomo} evaluates very-long-term \emph{two-party}
dialogue grounded in speaker personas and temporally ordered event graphs
(Table~\ref{tab:bench-comparison}). Its temporal-QA category touches evolving facts,
but supersession with explicit attribution is not a labeled task, and the corpus is
a single shared dyadic narrative rather than a multi-author organization. Top
systems now cluster between roughly 83\% and 93\%; tellingly, the numbers are also
\emph{disputed}, with different evaluators reporting materially different scores for
the same systems, which is itself an argument for a neutral, fixed-harness
benchmark.

\textbf{LongMemEval}~\citep{wu2024longmemeval} measures five memory abilities over
scalable user-assistant chat. Crucially, its knowledge-update and temporal
categories operate \emph{within one user's own stated facts}: there is no
multi-author corpus, no channel heterogeneity, and no role structure. Leading
systems now exceed 83\%.

\textbf{BEAM}~\citep{beam2025} pushes the \emph{volume} frontier, with ten memory
tasks (including knowledge update, contradiction resolution, event ordering, and
abstention) over coherent \emph{single-narrator} conversations generated up to 10M
tokens across two tracks (1M and 10M). At extreme scale it is not saturated:
reported best accuracy is about 64\% at 1M tokens and about 49\% at 10M. But the
challenge it poses is one of how \emph{much} a single narrator said, not of
\emph{who} said what across an organization; it varies length, not source
heterogeneity.

\textbf{MemoryAgentBench}~\citep{hu2025memoryagentbench} shifts toward agentic,
incremental multi-turn interaction over chunked inputs, a meaningful move beyond
static QA, but it still models one agent processing one input stream. Related
continual-learning and memory-hallucination benchmarks, including
MemoryBench~\citep{memorybench2025}, share the same single-user, personal-scale
framing and are not analyzed individually here.

The common thread is that near-saturation on LoCoMo and LongMemEval does not show
that memory is solved; it shows the benchmarks have stopped measuring the hard
part. Table~\ref{tab:bench-comparison} makes the gap concrete on the decisive
axes: \emph{modality}, where prior work is two-party, single-narrator, or
user--assistant dialogue while \bench{} is a multi-channel organizational corpus
authored by many distinct people; and \emph{scale}, where prior work grows one
narrator's history in \emph{length} (to ${\sim}1$M and 10M tokens) while \bench{}
instead grows \emph{breadth}---more authors, channels, and a multi-year timeline.
The remaining axes are where prior work is uniformly weak: none varies author
count (\emph{multi-source}), and although several include knowledge-update QA, none
treats \emph{supersession}---a fact invalidated and replaced, carrying an
attributed reason and date---as a first-class, labeled reasoning category.

\subsection{Long-Horizon and Temporal Reasoning}
\label{sec:related:temporal}

Beyond memory systems specifically, \bench{} draws on two strands of work on
temporal reasoning in language. The first concerns \emph{long-context} models and
retrieval-augmented generation, where a recurring finding is that the ability to
\emph{attend} to a long input is not the ability to \emph{reason} over information
distributed across it. The second concerns \emph{temporal} structure:
distinguishing valid-time (when a fact held in the world) from ingest-time (when it
was recorded), and reconstructing knowledge as of a past date. \bench{} makes this
structure first-class through the bi-temporal source-of-truth graph described in
\S\ref{sec:design}, and it is the dimension on which we observe the sharpest
degradation across all evaluated systems.

% COMPARISON TABLE: OrgMemBench vs existing benchmarks on the axes that define
% organizational memory. booktabs, no vertical rules (guide 1.6). SOTA / scale
% facts from landscape.md; OrgMemBench medium leaderboard is planned (see note).
\begin{table}[t]
  \centering
  \caption{\textbf{\bench{} versus existing memory benchmarks.} To our knowledge
    the first benchmark to target multi-source, multi-author organizational memory
    with bi-temporal supersession as a labeled reasoning category. The axes are
    walked through in the text (\S\ref{sec:related:benchmarks}).}
  \label{tab:bench-comparison}
  \small
  \setlength{\tabcolsep}{4pt}
  \begin{tabular}{@{}l p{2.5cm} p{2.5cm} l c c c@{}}
    \toprule
    & & & Temporal & Multi- & Super- & Public \\
    Benchmark & Modality & Scale & span & source & session & leaderboard \\
    \midrule
    LoCoMo
      & Dyadic dialogue + image
      & $\sim$9K tok/conv (50)
      & Up to 35 sess.
      & No
      & Partial
      & Informal \\
    LongMemEval
      & User--assistant chat
      & 128K\,$\rightarrow\sim$1M tok
      & Multi-session
      & No
      & Yes
      & Yes \\
    BEAM
      & Single-narrator dialogue
      & Up to 10M tok/conv
      & $\sim$1 year
      & No
      & Yes
      & Yes \\
    \midrule
    \bench{} (ours)
      & Multi-channel org.\ corpus
      & $\sim$63K (S) / ${>}200$K (M) tok
      & Multi-year
      & Yes
      & Yes
      & Planned \\
    \bottomrule
  \end{tabular}
  \par\smallskip
  {\footnotesize ``Partial'' for LoCoMo reflects its temporal-QA category, which
  touches evolving facts but does not label supersession with attribution. Token
  scales are approximate; \bench{} figures are plain-text corpus totals.}
\end{table}
